% Generated draft: requires mathematical and source review.
\documentclass[11pt]{article}
\usepackage[margin=25mm]{geometry}
\usepackage{fontspec}
\setmainfont{Latin Modern Roman}
\usepackage{amsmath,amssymb,mathtools,unicode-math}
\setmathfont{Latin Modern Math}
\usepackage{enumitem,needspace,xurl,hyperref}
\setlength{\parindent}{0pt}
\setlength{\parskip}{5pt}
\setlength{\emergencystretch}{3em}
\newcommand{\N}{\mathbb{N}}
\newcommand{\Z}{\mathbb{Z}}
\newcommand{\Q}{\mathbb{Q}}
\newcommand{\R}{\mathbb{R}}
\newcommand{\C}{\mathbb{C}}
\newcommand{\F}{\mathbb{F}}
\newcommand{\A}{\mathbb{A}}
\newcommand{\PP}{\mathbb{P}}
\newcommand{\E}{\mathbb{E}}
\newcommand{\eps}{\varepsilon}
\newcommand{\dd}{\,\mathrm{d}}
\begin{document}
\section*{UnsolvedMath ID: 30006073}
\textbf{Source problem code:} catalogue entry
\subsection{Definitions and mathematical statement}
Let \(p\) be a prime, let \(k/\Q_p\) be a finite extension, and fix an algebraic closure \(\bar k\). Let \(A/k\) be an abelian variety and let \(i:X\hookrightarrow A\) be a closed embedding over \(k\), where \(X/k\) is a smooth projective curve of genus at least \(2\).

Define the prime-to-\(p\) geometric torsion subgroup by
\[
A(\bar k)_{\mathrm{tors},p'}
=\{a\in A(\bar k):[m]a=0\text{ for some integer }m\geq 1\text{ with }\gcd(m,p)=1\}.
\]
Read the source hypothesis that \(X\) contains no torsion points of order prime to \(p\) as
\[
i\bigl(X(\bar k)\bigr)\cap A(\bar k)_{\mathrm{tors},p'}=\varnothing.
\]
This uses the geometric interpretation of the torsion hypothesis; the supplied excerpt does not explicitly distinguish geometric torsion points from \(k\)-rational torsion points.

For every integer \(n\geq 0\), let \([p^n]:A\to A\) denote multiplication by \(p^n\), and define the scheme-theoretic inverse image
\[
X_n=A\mathop{\times}_{A}X,
\]
where the maps to the common base \(A\) are \([p^n]\) on the first factor and \(i\) on the second factor. For a nonempty proper scheme \(Y/k\), define its index by
\[
\operatorname{ind}_k(Y)
=\gcd\{[\kappa(y):k]:y\in Y\text{ is a closed point}\},
\]
where \(\kappa(y)\) is the residue field of \(y\). This definition does not require \(X_n\) to be connected.

\textbf{Conjecture 1 (supplied source).} Under the preceding hypotheses, does there exist an integer \(N\geq 0\), allowed to depend on the given data, such that
\[
\forall n\in\Z_{\geq 0},\qquad n\geq N\ \Longrightarrow\ p\mid\operatorname{ind}_k(X_n)?
\]
Equivalently, for every \(n\geq N\), every closed point \(x\in X_n\) must satisfy \(p\mid[\kappa(x):k]\).

\textbf{Source-reported partial result and qualifications.} The excerpt reports that for every \(r>0\), there exists a sufficiently large integer \(n_r\) such that every closed point of \(X_{n_r}\) having ramification index at most \(r\) has degree divisible by \(p\). Under the convention that ramification index is measured relative to \(k\), this reads
\[
\forall r>0\ \exists n_r\geq 0\ \forall x\in X_{n_r}\text{ closed},\qquad
 e(\kappa(x)/k)\leq r\ \Longrightarrow\ p\mid[\kappa(x):k].
\]
The excerpt does not specify the reference field for the ramification index; this displayed formulation records the relative convention provisionally. The reported bound concerns points with bounded ramification and is not asserted here to establish the conjecture.

The excerpt also states consequences conditional on Conjecture 1 concerning toric Galois sections over finite extensions of \(\Q_p\), Selmer Galois sections over number fields, and standard Galois sections. Their notation is corrupted in the supplied transcription, and the statement labelled Theorem 3 is truncated after the introduction of a finite \(\acute{e}\)tale cover \(Y\to X\). These consequences are retained as conditional source context; their full mathematical statements cannot be reconstructed reliably from the supplied text.
\subsection{Short English statement}
For a curve of genus at least \(2\) embedded in an abelian variety over a finite extension of \(\Q_p\), assume that the curve contains no prime-to-\(p\) torsion point, interpreted geometrically. The conjecture asks whether every closed point of every sufficiently deep inverse image under multiplication by \(p^n\) has degree divisible by \(p\). The supplied source reports a partial result for points of bounded ramification and makes its subsequent section results conditional on this conjecture. This is a statement-only draft requiring mathematical review.
\subsection{Sources}
\begin{itemize}
\item Local registry: \texttt{data/unsolvedmath/problems.json}, permanent UnsolvedMath ID \(30006073\), source problem code \texttt{catalogue entry}, title \emph{Index Divisibility in p-Power Pullbacks of Curves}, in the supplied collection \emph{Oberwolfach Reports Open Problems}.
\item Supplied report citation: \emph{Anabelian Geometry and Representations of Fundamental Groups} (2025), Oberwolfach Reports, DOI \texttt{10.4171/owr/2024/45}, pp.\ 2621--2670. \url{https://doi.org/10.4171/owr/2024/45}. The supplied excerpt labels the target assertion Conjecture 1 and ends with a page marker for p.\ 2635. No individual proposer is identified in the record.
\item Supplied literature source: \emph{Oberwolfach Report 45/2024}. \url{https://doi.org/10.4171/owr/2024/45}.
\item Literature assessment dated exactly \texttt{2026-08-22}: the supplied curation labels the conjecture open, reports that no proof or counterexample was found in its primary-literature check, and says that downstream section results remain conditional on the conjecture. This is supplied dated curation, not an independently verified literature assessment.
\item The registry marks the statement \texttt{corrected\_verified} and reports that the official report was checked and its notation repaired. These are supplied provenance claims, not verification performed for this draft. The citation's year 2025 and the report identifier 45/2024 are preserved as supplied.
\end{itemize}
\end{document}
